\section{Results}

Table~\ref{tab:singular-values} reports the singular values of the scaled
Jacobian for the five-mode map
$\mathbf{f}=[f_2,\ldots,f_6]^\mathsf{T}$, using central differences with
relative step $\Delta\theta_j/\theta_j=\num{1e-4}$.  The equivalent-sphere row
is the baseline: because that model depends on $a$ and $c$ only through
$R_\mathrm{eq}$, the exact Jacobian is rank deficient and the reported
$\sigma_3=\num{1.54e-10}$ is only the floating-point residue of a zero singular
value.  The Ritz sphere is therefore a distinct calculation.  It retains the
same parameter vector and mode set, but evaluates the full spheroidal Ritz
model at $a=c$.

The oblate branch then shows the relevant lifting behaviour.  Moving from the
Ritz sphere to the near-oblate point at $\varepsilon=\num{0.436}$
($c/a=\num{0.9}$) raises the
weakest singular value from $\sigma_3=\num{0.00806}$ to
$\sigma_3=\num{0.01274}$, and the canonical oblate point
$(c/a=2/3,\ \varepsilon\approx\num{0.745})$ raises it further to
$\sigma_3=\num{0.02643}$.  The corresponding condition number falls from
$\kappa_\mathrm{Ritz}=244$ to $\kappa=145$ and then to
$\kappa_\mathrm{oblate}=\num{69.4}$.  Since the inverse variables and measured
mode set are unchanged, that improvement is geometric rather than procedural.

The prolate continuation does not show systematic improvement.  Along this
branch $a=a_0$ is held fixed while $c$ increases (see Section~2).  Across the tested
range $c/a\in[\num{1.0},\num{1.5}]$, $\kappa_\mathrm{prolate}$ remains in the
range 156--244, showing no monotonic decrease with increasing asphericity.  At
the representative endpoint $c/a=\num{1.5}$ the scaled Jacobian gives
$\sigma_3=\num{0.01636}$ and $\kappa=156$, still markedly worse than the
canonical oblate configuration ($\kappa=\num{69.4}$).

\begin{table}[t]
  \centering
  \small
  \caption{Singular values of the scaled Jacobian
  $\widetilde{J}=\operatorname{diag}(\mathbf{f})^{-1}J
  \operatorname{diag}(\boldsymbol{\theta})$ for the five-mode inverse map
  $\mathbf{f}=[f_2,\ldots,f_6]^\mathsf{T}$.  Central differences use
  $\Delta\theta_j/\theta_j=\num{1e-4}$.  The Ritz-sphere $\sigma_3$ is a
  discretisation artefact; $\sigma_3=0$ in the continuous limit (see text).
  $\varepsilon$ denotes oblate eccentricity; the prolate endpoint is
  parameterised by $c/a$ as defined in Section~2.}
  \label{tab:singular-values}
  \begin{tabular}{@{}lcccccc@{}}
    \toprule
    Case & $c/a$ & $\varepsilon$ & $\sigma_1$ & $\sigma_2$ & $\sigma_3$ & $\kappa$ \\
    \midrule
    Equivalent sphere & \num{1.0} & \num{0} &
    \num{1.893} & \num{0.112} & \num{1.54e-10} & \num{1.23e10} \\
    Ritz sphere & \num{1.0} & \num{0} &
    \num{1.971} & \num{0.573} & \num{0.00806} & \num{244} \\
    Near oblate & \num{0.9} & \num{0.436} &
    \num{1.854} & \num{0.436} & \num{0.01274} & \num{145} \\
    Canonical oblate & $2/3$ & \num{0.745} &
    \num{1.835} & \num{0.130} & \num{0.02643} & \num{69.4} \\
    Prolate endpoint & \num{1.5} & --- &
    \num{2.551} & \num{0.382} & \num{0.01636} & \num{156} \\
    \bottomrule
  \end{tabular}
\end{table}

The near-sphere behaviour of the weakest singular value is governed by
\begin{equation}
  \sigma_3(\varepsilon)
  =
  \lambda_1 \varepsilon^2 + O(\varepsilon^4),
  \qquad \lambda_1>0.
  \label{eq:sigma3-expansion}
\end{equation}
At the exact sphere ($\varepsilon=0$), spherical symmetry requires
$\sigma_3=0$: the two geometric columns of $\widetilde{J}$ both reduce to
derivatives with respect to $R_\mathrm{eq}$ alone and are therefore
proportional.  A finite Ritz computation at $a=c$ does return a small positive
value ($\sigma_3=\num{0.00806}$ in Table~\ref{tab:singular-values}); however,
this is a discretisation artefact.  The present Ritz basis does not preserve the
exact $a\leftrightarrow c$ exchange symmetry of the sphere; accordingly, the
spherical rank degeneracy $\sigma_3(0)=0$ must be inferred from the continuous
model rather than from the discrete intercept.  Concretely, the
oblate-spheroidal integration metric and the fluid-mass blending coefficients
break that symmetry, and the resulting intercept \emph{grows} rather than
converges with basis size: $\sigma_0\approx\num{0.006}$, \num{0.011}, and
\num{0.018} for the three-, five-, and seven-mode Ritz bases respectively (here
the number refers to the Legendre-polynomial trial functions retained; see
Section~2).  The quadratic onset $\sigma_3\sim\lambda_1\varepsilon^2$ is a
symmetry-based expectation consistent with the oblate trend.  In the present
discrete implementation, the artefactual intercept remains comparable to the
quadratic term for $\varepsilon\lesssim\num{0.5}$, so the clean asymptotic
regime is the canonical operating range ($\varepsilon\approx\num{0.75}$) rather
than the near-sphere neighbourhood.

A least-squares fit of $\sigma_3(\varepsilon)=\sigma_0+\lambda_1\varepsilon^2$ to the
oblate data yields $\sigma_0=\num{0.0074}$ and $\lambda_1=\num{0.034}$; the
non-zero intercept in this fit is the same discretisation effect described
above, not a physical floor.

The scientific conclusion therefore rests on the deformation-dependent quadratic
growth in \eqref{eq:sigma3-expansion}.  At the canonical oblate operating point
($\varepsilon\approx\num{0.745}$), the quadratic term is well above the
artefactual intercept, so the oblate branch separates the weak inverse direction
more effectively than either the equivalent-sphere model or the prolate
continuation.  The reported condition numbers at the canonical oblate point
($\kappa\approx 69$) are computed directly and are not affected by the
near-sphere artefact.

The practical meaning is straightforward.  The second flexural mode
$f_2=\SI{3.95}{\hertz}$ already lies in the low-frequency band used by the
companion studies, and the higher modes in the set respond differently to the
same curvature perturbation.  That differential response separates the
sensitivity directions for $a$, $c$, and $E$ in the oblate family.  The
prolate family does not deliver the same separation, so loss of spherical
symmetry is not, by itself, a useful identifiability category.
